\def\tituloconcepto{Control cuántico}
% Si el capitulo se compila solo (sin main.tex), esta cabecera abre el documento.
% Antes de % Si el capitulo se compila solo (sin main.tex), esta cabecera abre el documento.
% Antes de % Si el capitulo se compila solo (sin main.tex), esta cabecera abre el documento.
% Antes de \input{cabecera} cada capitulo define \tituloconcepto.
\ifdefined\enmain\else
  \documentclass[11pt,a4paper]{article}
  \input{preambulo}
  \begin{document}
  \title{\tituloconcepto}
  \author{\autor}
  \date{\today}
  \maketitle
\fi
 cada capitulo define \tituloconcepto.
\ifdefined\enmain\else
  \documentclass[11pt,a4paper]{article}
  % Preambulo compartido por todos los documentos. Compilar con XeLaTeX/LuaLaTeX
% (o Tectonic). Si se usa Overleaf, elegir el compilador XeLaTeX.
\usepackage[spanish,es-noshorthands]{babel}
\usepackage{amsmath,amssymb}
\usepackage{booktabs,array}
\usepackage[a4paper,margin=2.5cm]{geometry}
\usepackage{xcolor}
\usepackage{listings}
\usepackage[most]{tcolorbox}
\usepackage{quantikz}
\usepackage[hidelinks]{hyperref}

% ---- Autor (editar a gusto) -------------------------------------------------
\newcommand{\autor}{Genesis Silesky}

% ---- Notacion de Dirac y atajos ---------------------------------------------
% quantikz ya define \ket y \bra; DeclareRobustCommand los redefine sin error.
\DeclareRobustCommand{\ket}[1]{\left|#1\right\rangle}
\DeclareRobustCommand{\bra}[1]{\left\langle#1\right|}
\DeclareRobustCommand{\braket}[2]{\left\langle#1\middle|#2\right\rangle}
\newcommand{\C}{\mathbb{C}}
\newcommand{\R}{\mathbb{R}}
\DeclareMathOperator{\Tr}{Tr}

% ---- Cajas de apoyo --------------------------------------------------------
% "concepto": definiciones previas que hacen falta para entender el capitulo.
\newtcolorbox{concepto}[1]{breakable,colback=blue!4,colframe=blue!55!black,
  fonttitle=\bfseries,title={Concepto previo: #1}}
% "nota": aclaraciones y advertencias.
\newtcolorbox{nota}[1]{breakable,colback=orange!6,colframe=orange!70!black,
  fonttitle=\bfseries,title={#1}}

% ---- Estilos de listados ----------------------------------------------------
\lstdefinestyle{py}{language=Python,basicstyle=\ttfamily\footnotesize,
  breaklines=true,postbreak=\mbox{\textcolor{gray}{$\hookrightarrow$}\space},
  numbers=left,numberstyle=\tiny\color{gray},frame=single,
  commentstyle=\color{green!40!black}\itshape,
  keywordstyle=\color{blue!70!black}\bfseries,
  stringstyle=\color{red!60!black},showstringspaces=false,tabsize=4,
  columns=fullflexible,keepspaces=true}
\lstdefinestyle{salida}{basicstyle=\ttfamily\footnotesize,breaklines=true,
  frame=leftline,backgroundcolor=\color{gray!8},columns=fullflexible,
  keepspaces=true}

  \begin{document}
  \title{\tituloconcepto}
  \author{\autor}
  \date{\today}
  \maketitle
\fi
 cada capitulo define \tituloconcepto.
\ifdefined\enmain\else
  \documentclass[11pt,a4paper]{article}
  % Preambulo compartido por todos los documentos. Compilar con XeLaTeX/LuaLaTeX
% (o Tectonic). Si se usa Overleaf, elegir el compilador XeLaTeX.
\usepackage[spanish,es-noshorthands]{babel}
\usepackage{amsmath,amssymb}
\usepackage{booktabs,array}
\usepackage[a4paper,margin=2.5cm]{geometry}
\usepackage{xcolor}
\usepackage{listings}
\usepackage[most]{tcolorbox}
\usepackage{quantikz}
\usepackage[hidelinks]{hyperref}

% ---- Autor (editar a gusto) -------------------------------------------------
\newcommand{\autor}{Genesis Silesky}

% ---- Notacion de Dirac y atajos ---------------------------------------------
% quantikz ya define \ket y \bra; DeclareRobustCommand los redefine sin error.
\DeclareRobustCommand{\ket}[1]{\left|#1\right\rangle}
\DeclareRobustCommand{\bra}[1]{\left\langle#1\right|}
\DeclareRobustCommand{\braket}[2]{\left\langle#1\middle|#2\right\rangle}
\newcommand{\C}{\mathbb{C}}
\newcommand{\R}{\mathbb{R}}
\DeclareMathOperator{\Tr}{Tr}

% ---- Cajas de apoyo --------------------------------------------------------
% "concepto": definiciones previas que hacen falta para entender el capitulo.
\newtcolorbox{concepto}[1]{breakable,colback=blue!4,colframe=blue!55!black,
  fonttitle=\bfseries,title={Concepto previo: #1}}
% "nota": aclaraciones y advertencias.
\newtcolorbox{nota}[1]{breakable,colback=orange!6,colframe=orange!70!black,
  fonttitle=\bfseries,title={#1}}

% ---- Estilos de listados ----------------------------------------------------
\lstdefinestyle{py}{language=Python,basicstyle=\ttfamily\footnotesize,
  breaklines=true,postbreak=\mbox{\textcolor{gray}{$\hookrightarrow$}\space},
  numbers=left,numberstyle=\tiny\color{gray},frame=single,
  commentstyle=\color{green!40!black}\itshape,
  keywordstyle=\color{blue!70!black}\bfseries,
  stringstyle=\color{red!60!black},showstringspaces=false,tabsize=4,
  columns=fullflexible,keepspaces=true}
\lstdefinestyle{salida}{basicstyle=\ttfamily\footnotesize,breaklines=true,
  frame=leftline,backgroundcolor=\color{gray!8},columns=fullflexible,
  keepspaces=true}

  \begin{document}
  \title{\tituloconcepto}
  \author{\autor}
  \date{\today}
  \maketitle
\fi


\section{Control cuántico: Hamiltoniano, rotaciones y acoplamiento}

\subsection{Objetivo}
Entender de dónde salen las puertas cuánticas: cada una es una evolución
unitaria $U(t)=e^{-iHt}$ que se consigue eligiendo el Hamiltoniano $H$.
Se ve cómo se hacen rotaciones de un qubit, qué pasa con un pulso desafinado y
cómo el acoplamiento entre dos qubits los entrelaza.

\subsection{Conceptos previos}
\begin{concepto}{unitariedad, matrices de Pauli y producto tensorial}
Ver sección~\ref{sec:prerreq} (matrices unitarias y producto tensorial) y los
capítulos~10 y~11. Las matrices de Pauli son
$\sigma_x=\begin{psmallmatrix}0&1\\1&0\end{psmallmatrix}$,
$\sigma_y=\begin{psmallmatrix}0&-i\\i&0\end{psmallmatrix}$ y
$\sigma_z=\begin{psmallmatrix}1&0\\0&-1\end{psmallmatrix}$.
\end{concepto}

\subsection{Fundamento teórico}
\paragraph{Unitariedad.} Una operación $U$ es unitaria si $U^\dagger U=I$. Los
postulados de la mecánica cuántica dicen que toda evolución física, incluidas
las operaciones lógicas de un computador cuántico, es unitaria. Su significado
físico es que la suma de probabilidades se conserva:
\begin{equation}
  \langle U\psi|U\psi\rangle=\langle\psi|U^\dagger U|\psi\rangle
  =\langle\psi|\psi\rangle=|\alpha|^2+|\beta|^2 .
\end{equation}
No se pierde ni se crea nada: la operación no destruye información.

\paragraph{Hamiltoniano y evolución.} El Hamiltoniano $H$ es el operador de la
energía total del sistema y manda su evolución en el tiempo (ecuación de
Schrödinger):
\begin{equation}
  i\hbar\frac{d}{dt}\ket{\psi(t)}=H\ket{\psi(t)} .
\end{equation}
Con $\hbar=1$ y un $H$ constante, la solución es $\ket{\psi(t)}=U(t)\ket{\psi(0)}$
con
\begin{equation}
  U(t)=e^{-iHt} .
\end{equation}
\emph{Controlar} un computador cuántico es, básicamente, modificar $H$ para
obtener la evolución $U(t)$ deseada. En general $H=H_0+H_1$: la parte
\emph{interna} $H_0$ (energía Zeeman de cada qubit y acoplamiento entre ellos)
y la parte \emph{externa} $H_1$ (el campo de radiofrecuencia aplicado).

\paragraph{Un qubit: rotaciones.} Un pulso resonante de duración $t$, frecuencia
de Rabi $\omega_1$ y fase $\varphi$ produce
\begin{equation}
  U=\exp\!\Big[-i\,\pi\omega_1t\,\big(\cos\varphi\,\sigma_x+\sin\varphi\,\sigma_y\big)\Big],
\end{equation}
una \textbf{rotación} alrededor del eje $\cos\varphi\,\hat x+\sin\varphi\,\hat y$ de
la esfera de Bloch, por un ángulo $\theta=2\pi\omega_1t$. La fase del pulso
elige el eje y la duración elige el ángulo. Se escribe $R_x(\theta)$,
$R_y(\theta)$, etc.

\paragraph{La rotación alrededor de $z$.} Un solo pulso no puede girar
alrededor de $z$, pero se puede armar con tres:
\begin{enumerate}
  \item girar $\pi/2$ alrededor de $y$ (el eje $x$ ocupa el lugar de $z$);
  \item girar $\theta$ alrededor de $x$ (equivale a girar alrededor de $z$);
  \item girar $\pi/2$ alrededor de $-y$ para volver al sistema original.
\end{enumerate}
Así, con pulsos de radiofrecuencia se puede girar un qubit alrededor de
cualquier eje. El \textbf{orden de los pulsos importa}, porque las
rotaciones, en general, no conmutan.

\paragraph{Pulso desafinado.} Si la frecuencia del pulso difiere de la del
qubit en $\Delta\omega$, la evolución combina dos giros a la vez: alrededor de
$z$ (a frecuencia $\Delta\omega$) y alrededor del eje del pulso. Con una
desintonía grande el pulso no logra excitar al qubit.

\paragraph{Dos qubits acoplados.} Con acoplamiento $J_{12}$, la parte interna
del Hamiltoniano (en el sistema doblemente rotante) es
$H_0=2\pi J_{12}\,\sigma_z\otimes\sigma_z/4$. Es una rotación alrededor del eje
$\sigma_z\otimes\sigma_z$, que equivale a decir que \emph{el eje de giro de un
qubit depende del estado del otro}: si el otro está en $\ket0$ el eje es $+z$ y
si está en $\ket1$ es $-z$. Si el otro qubit está en una superposición, la
frecuencia de precesión también queda en superposición y los dos qubits se
\textbf{entrelazan}.

\paragraph{Criterio de entrelazamiento.} Si los qubits no están entrelazados,
el estado total es un producto $\rho=\rho_A\otimes\rho_B$. Si $\rho$ no se
puede escribir así, el sistema está entrelazado.

\subsection{Código (completamente comentado)}
\lstinputlisting[style=py]{../../fuentes/12_control_cuantico.py}

\subsection{Explicación celda por celda}
\begin{description}
  \item[Celda 2] Calcula $U=e^{-iHt}$ para un $H$ cualquiera y comprueba
    $U^\dagger U=I$ y que la suma de probabilidades sigue valiendo 1.
  \item[Celda 3] Define $R_x$, $R_y$, $R_z$ y verifica que un pulso de
    $0.25$~ms a $1000$~Hz es $R_x(\pi/2)$.
  \item[Celda 4] Arma $R_z(\theta)$ con la secuencia de tres pulsos y
    comprueba que el orden contrario da $R_z(-\theta)$.
  \item[Celda 5] Aplica un pulso $\pi$ con distintas desintonías.
  \item[Celda 6] Con $J=696$~Hz y $t=1/(2J)$ deja al qubit 0 en $\ket+$ y
    muestra que su fase cambia de signo según el estado del qubit 1.
  \item[Celda 7] Comprueba, desde $\ket{++}$, que tras $t=1/(2J)$ el estado
    ya no es un producto.
\end{description}

\subsection{Salida obtenida}
\lstinputlisting[style=salida]{../../fuentes/salidas/12_control_cuantico.txt}
Es determinista (no hay medición).

\subsection{Interpretación}
La suma de probabilidades queda en 1 antes y después de $U$. Los tres pulsos
$R_y(\pi/2)$, $R_x(\theta)$, $R_y(-\pi/2)$ dan $R_z(+\theta)$, y con el orden
cambiado se obtiene $R_z(-\theta)$. Con desintonía de $1000$~Hz un pulso $\pi$
solo logra $P(1)=0.317$, y con $10\,000$~Hz, $0.000$. Con el qubit 1 en
$\ket0$ o en $\ket1$, el qubit 0 acaba en $(x,y)=(0,+1)$ o $(0,-1)$: fases
opuestas. Tras $t=0.718$~ms el estado ya no es un producto y la entropía de un
qubit es $1$ bit (máximo).

\subsection{Ejercicios}
\begin{enumerate}
  \item Construye $R_z(\theta)$ con el orden contrario de pulsos
    ($R_y(-\pi/2)$ primero) y comprueba qué signo obtienes.
  \item ¿Qué tiempo $t$ hace falta con $J=696$~Hz para una rotación de
    $\pi$ en el eje $\sigma_z\otimes\sigma_z$?
  \item En la Celda~7, calcula la entropía para varios tiempos entre $0$ y
    $1/J$ y observa cuándo es máxima.
\end{enumerate}

\subsection{Errores comunes}
\begin{itemize}
  \item Multiplicar las matrices en el orden en que se aplican los pulsos:
    en un producto de matrices el primer pulso aplicado queda a la
    \emph{derecha}.
  \item Con \texttt{np.kron(A, B)}, el factor de la izquierda es el qubit 1 y el
    de la derecha el qubit 0 (orden de Qiskit).
\end{itemize}

\ifdefined\enmain\else
  \end{document}
\fi

